\subsection{Text \& Image Retrieval}
\label{sec:innov-retrieval}

The second innovation is a retrieval stage that draws on \emph{both} the enriched
text and the (text / image) vectors of Section~\ref{sec:innov-offline}, and that is
composed adaptively according to what the user actually provided. Two workers ---
Text Retrieval and Visual Retrieval --- run in parallel under a single Recall
agent; within Text Retrieval, a \emph{keyword} channel and a \emph{text-semantic}
channel are fused first, after which the text and visual channels are fused by
Reciprocal Rank Fusion (RRF).

\subsubsection{Text retrieval}
\label{sec:text}

The user request is first turned into a preference profile by an upstream
analysis LLM, whose English \emph{search keywords} $K_q$ are the primary query
signal. The Text Retrieval worker forms a query string from $K_q$ (falling back
to a planner hint or a preference-derived string when $K_q$ is empty), applies a
budget-based price ceiling and a platform post-filter, and runs two
complementary channels whose ranked lists are later RRF-merged.

\paragraph{Keyword channel.}
The query is tokenized into a content token set $Q$ (length $\geq 3$, stop-words
removed, de-duplicated, capped). A product is recalled if any token matches its
product name or enriched text:
\begin{equation}
  \mathrm{recall}(p_i)=\bigvee_{k\in Q}
  \big(k \sqsubseteq \mathrm{name}(p_i)\ \lor\ k \sqsubseteq t_i\big),
  \label{eq:text-recall}
\end{equation}
where $\sqsubseteq$ denotes substring containment. Matching against $t_i$ is the
payoff of the enrichment layer: fine-grained visual and experiential terms become
searchable even though they are absent from the product name. Candidates are then
ranked by a name-plus-enriched relevance score
\begin{equation}
  \mathrm{rel}(p_i)=\sum_{k\in Q}\Big(
    \mathbb{1}\!\left[k\sqsubseteq\mathrm{name}(p_i)\right]
    +\mathbb{1}\!\left[k\sqsubseteq t_i\right]
  \Big),
  \label{eq:rel}
\end{equation}
with ties broken by popularity, rating or price. Category is intentionally
\emph{not} a hard filter here: recall stays high and category precision is
delegated to the Reject/Verify agent (Section~\ref{sec:innov-reject}).

\paragraph{Text-semantic channel.}
In parallel, the same worker embeds the query string with the shared multimodal
encoder $\phi$ of Section~\ref{sec:innov-offline} (DashScope
\texttt{multimodal-embedding-v1} by default; overridable via
\texttt{QWEN\_MM\_EMBED\_MODEL}) and retrieves by cosine similarity against
offline product text vectors $\mathbf{u}_i$ stored in
\texttt{text\_embedding}. Each product contributes \emph{one} vector: the embed
text is $t_i$ when present, otherwise a concatenation of name, summary and
category, truncated to $L{=}2000$ characters --- there is no multi-chunk
splitting. The channel returns
\begin{equation}
  s^{\mathrm{txt}}_i=\cos\!\big(\phi(q),\,\mathbf{u}_i\big),
  \qquad
  \mathrm{TextSemRecall}(q)=\operatorname*{top\text{-}K}_{i}\, s^{\mathrm{txt}}_i,
  \label{eq:stxt}
\end{equation}
optionally restricted by a soft category guardrail when a target category is
known. Each hit is annotated with \texttt{\_text\_score}. The channel is gated by
\texttt{VS\_TEXT\_SEMANTIC} and activates only when the catalog contains text
vectors; otherwise Text Retrieval degrades gracefully to the keyword channel
alone.

\paragraph{Intra-text RRF.}
Let $L^{\mathrm{kw}}$ and $L^{\mathrm{sem}}$ be the two ranked lists (1-based
ranks). They are fused by Reciprocal Rank Fusion with constant $k_{\mathrm{rrf}}$
(default $60$):
\begin{equation}
  \mathrm{RRF}(p)=\sum_{\ell\in\{\mathrm{kw},\mathrm{sem}\}}
  \frac{w_{\ell}}{k_{\mathrm{rrf}}+\mathrm{rank}_{\ell}(p)},
  \label{eq:rrf-text}
\end{equation}
where missing lists contribute nothing and $w_{\ell}$ defaults to $1$. The fused
list is the Text Retrieval output and carries \texttt{\_rrf\_score}.

\begin{chalgorithm}{Online text retrieval (keyword $\oplus$ semantic)}\label{alg:trecall}
  \item \textbf{Input:} preference profile (keywords $K_q$, budget, platform),
        optional query hint; text index $\{(\mathrm{id}_i,\mathbf{u}_i)\}$.
  \item Build query string $q$ from $K_q$ (else hint / preference fallback).
  \item $L^{\mathrm{kw}}\gets$ keyword SQL recall on $\mathrm{name}$ / $t_i$
        (Eq.~\ref{eq:text-recall}--\ref{eq:rel}), with price and platform filters.
  \item \textbf{if} text-semantic is enabled \textbf{and} the text index is
        non-empty \textbf{then}
        $L^{\mathrm{sem}}\gets\operatorname{top\text{-}K}\cos(\phi(q),\mathbf{u}_i)$
        \textbf{else} $L^{\mathrm{sem}}\gets\varnothing$.
  \item \textbf{return} $\mathrm{RRF}(L^{\mathrm{kw}},L^{\mathrm{sem}})$
        (Eq.~\ref{eq:rrf-text}), attaching \texttt{\_text\_score} /
        \texttt{\_rrf\_score} as available.
\end{chalgorithm}

\subsubsection{Visual retrieval}
\label{sec:photo}

The Visual Retrieval worker retrieves products whose image is visually similar
to a user photo. On first use it loads all stored vectors into an in-memory
\texttt{VisualIndex} (ids, vectors $\{\mathbf{v}_i\}$, and categories parsed from
the visual layer); the index is cached across requests and invalidated on
re-enrichment. Given a user photo $I_q$, every product is scored by cosine
similarity in the \emph{same} embedding space $\phi$ as text:
\begin{equation}
  s^{\mathrm{vis}}_i=\cos\!\big(\phi(I_q),\,\mathbf{v}_i\big),
  \qquad
  \cos(\mathbf{a},\mathbf{b})=\frac{\mathbf{a}^{\top}\mathbf{b}}{\lVert\mathbf{a}\rVert\,\lVert\mathbf{b}\rVert},
  \label{eq:svis}
\end{equation}
When a target category $c$ is known --- for example inferred by the VLM from the
photo in the image-only modality --- a \emph{category guardrail} restricts
scoring to compatible products, preventing visually-similar but
categorically-wrong matches:
\begin{equation}
  \mathrm{VisualRecall}(I_q)=\operatorname*{top\text{-}K}_{\,i\,:\,\mathrm{cat}_i\sim c}\ s^{\mathrm{vis}}_i .
  \label{eq:vrecall}
\end{equation}
Each returned product is annotated with its \texttt{\_visual\_score}
(Algorithm~\ref{alg:vrecall}).

\begin{chalgorithm}{Online visual recall}\label{alg:vrecall}
  \item \textbf{Input:} user photo $I_q$, optional category $c$, index
        $\{(\mathrm{id}_i,\mathbf{v}_i,\mathrm{cat}_i)\}$.
  \item \textbf{if} the index is empty \textbf{then return} $\varnothing$.
  \item $\mathbf{q}\gets\phi(I_q)$.
  \item $S \gets \{(\mathrm{id}_i,\ \cos(\mathbf{q},\mathbf{v}_i)) : c\text{ unset}\ \lor\ \mathrm{cat}_i\sim c\}$.
  \item Sort $S$ by score in descending order; \textbf{return} the top-$K$ items,
        attaching \texttt{\_visual\_score} to each.
\end{chalgorithm}

\subsubsection{Fusion and modality-adaptive routing}
\label{sec:fusion}

The Recall agent fuses the Text Retrieval list $L^{\mathrm{txt}}$ and the Visual
Retrieval list $L^{\mathrm{vis}}$ by a second, modality-weighted RRF,
\begin{equation}
  \mathrm{RRF}_{\mathrm{mod}}(p)=
  \frac{w_{\mathrm{txt}}}{k_{\mathrm{rrf}}+\mathrm{rank}_{\mathrm{txt}}(p)}
  +\frac{w_{\mathrm{vis}}}{k_{\mathrm{rrf}}+\mathrm{rank}_{\mathrm{vis}}(p)},
  \label{eq:rrf-mod}
\end{equation}
with default weights $(w_{\mathrm{txt}},w_{\mathrm{vis}})$ depending on the
detected modality (e.g.\ $(1,0)$ for text-only, $(0.55,0.45)$ for text+image,
$(0.35,0.65)$ for image-only). Products that appear in both lists receive a
higher RRF score; a channel that is absent or empty contributes nothing. After
fusion, the agent attaches a visual score and (when a text index exists) a text
score to \emph{every} surviving candidate --- including those recalled by only
one channel --- so the downstream Recommend stage sees uniform
$s^{\mathrm{vis}}$ / $s^{\mathrm{txt}}$ signals and can re-rank primarily by
\texttt{\_rrf\_score}.

Which channels run is decided per request by a modality router that classifies
the input as text-only ($T$), text-plus-image ($T{+}I$), or image-only ($I$):
$T$ runs Text Retrieval only (keyword $\oplus$ text-semantic); $T{+}I$ runs both
workers and RRF-merges; $I$ first infers a category from the photo and then runs
category-guarded visual retrieval (optionally still admitting text if the
router so decides). Every channel is independently switchable through an
experiment configuration (\texttt{VS\_TEXT\_SEMANTIC}, \texttt{VS\_VISUAL\_RECALL},
enrichment flags), so the contribution of enrichment, text-semantic recall and
visual recall can be measured by ablation. Visual retrieval is strictly
additive: it activates only when a query image is present \emph{and} the catalog
contains image vectors; text-semantic recall likewise requires stored text
vectors; otherwise the pipeline proceeds on the remaining active channels.
